% =====================================================================
%  CROWN-Inpaint: Updated Project Proposal
%  ---------------------------------------------------------------------
%  CS 754 Course Project, Spring 2026, IIT Bombay
%  Compiles standalone with `pdflatex` (uses only standard packages).
% =====================================================================
\documentclass[11pt,a4paper]{article}

\usepackage[margin=1in]{geometry}
\usepackage{amsmath,amssymb,amsthm,bm,mathtools}
\usepackage{enumitem}
\usepackage{booktabs}
\usepackage{array}
\usepackage{hyperref}
\usepackage{xcolor}
\usepackage{titlesec}
\usepackage{microtype}
\usepackage{cite}

\hypersetup{
  colorlinks=true,
  linkcolor=black!70!blue,
  citecolor=black!50!green,
  urlcolor=black!70!blue,
}

% --------- Theorem environments ---------
\newtheorem{theorem}{Theorem}
\newtheorem{proposition}[theorem]{Proposition}
\newtheorem{lemma}[theorem]{Lemma}
\newtheorem{corollary}[theorem]{Corollary}
\theoremstyle{definition}
\newtheorem{definition}{Definition}
\theoremstyle{remark}
\newtheorem*{remark}{Remark}

% --------- Macros ---------
\newcommand{\R}{\mathbb{R}}
\newcommand{\E}{\mathbb{E}}
\newcommand{\Loss}{\mathcal{L}}
\newcommand{\norm}[1]{\left\lVert #1 \right\rVert}
\newcommand{\diag}{\operatorname{diag}}
\newcommand{\KNN}{\operatorname{KNN}}
\newcommand{\sign}{\operatorname{sign}}
\newcommand{\dvec}{\bm{d}}
\newcommand{\hvec}{\bm{h}}
\newcommand{\xvec}{\bm{x}}
\newcommand{\alphavec}{\bm{\alpha}}
\newcommand{\Sk}{\mathcal{S}_k}

% --------- Title block ---------
\title{\textbf{CROWN-Inpaint}: A Confidence-weighted, Regime-aware, \\
       Non-local, Multi-scale Iterative Framework for Image Inpainting \\
       \vspace{0.3em}
       \large -- Updated Project Proposal --}

\author{%
  CS 754 Project Group \\
  \textit{Department of Electrical Engineering, IIT Bombay} \\
  \texttt{\{23b0313, …\}@iitb.ac.in}}

\date{April 2026}

\begin{document}
\maketitle

\begin{abstract}
\noindent
We propose \textbf{CROWN-Inpaint}, a unified iterative framework for
single-image inpainting that integrates four complementary priors --
\emph{Confidence}, \emph{Regime}, \emph{Non-local}, and multi-scale
\emph{Warm-start} -- into a single optimisation-first solver layered on
top of an established masked dictionary-learning backbone. After a
careful audit of prior art (in particular Mairal--Elad--Sapiro 2008,
Naumova--Schnass 2018, and the recent Innsbruck Bachelor thesis of
Tiefenthaler 2018, which our group identified during literature
review), we explicitly position \emph{masked K-SVD} as a settled
baseline rather than a contribution, and re-centre our novelty on the
\textbf{coupling layer} above it. Our four novel mechanisms are: (i) a
continuous, time-varying confidence map that drives a weighted OMP
generalising the binary masked OMP of prior work; (ii) a per-pixel
regime map that fuses a biharmonic smooth-prior branch with the sparse
branch; (iii) a non-local centralised $\ell_1$ coefficient-coupling
step adapted from CSR/GSR but \emph{embedded inside} the iterative
masked solver, with a slow-rate Lasso oracle inequality bounding its
error by a Self-Similarity Index; and (iv) a coarse-to-fine pyramid
with prior injection that provably removes the
\emph{zero-information plateau} on which single-scale K-SVD gets stuck
when the hole exceeds the patch size. We provide a self-contained
convergence and error analysis, an ablation-rich experimental plan,
and a clear delineation of which components are baselines and which
are our actual contribution.
\end{abstract}

% =====================================================================
\section{Problem Statement}
% =====================================================================

Let $y\in\R^{H\times W}$ be a single-channel image observed through a
binary mask $M\in\{0,1\}^{H\times W}$, with $M_i=1$ iff pixel $i$ is
observed and $M_i=0$ iff pixel $i$ is missing. The observed set is
$\Omega=\{i:M_i=1\}$ and the hole is $\Omega^c$. We seek a
reconstruction $u\in\R^{H\times W}$ such that
\begin{enumerate}[label=(\roman*),itemsep=2pt]
\item $u_i = y_i$ for all $i\in\Omega$ (\emph{exact preservation of
  observed data}); and
\item $u$ on $\Omega^c$ is statistically and visually consistent with
  the observed surround.
\end{enumerate}
Within the dictionary-based paradigm, every $p\times p$ patch
$R_q u\in\R^n$ ($n=p^2$) is modelled as
\begin{equation}
R_q u \;\approx\; D \alphavec_q,\qquad \norm{\alphavec_q}_0\le s,
\label{eq:patch-model}
\end{equation}
for a learned dictionary $D\in\R^{n\times K}$ ($K\gtrsim n$,
unit-norm columns) and a sparse code
$\alphavec_q\in\R^K$ with at most $s$ non-zeros. Because the mask
forbids access to ground-truth patches, the dictionary itself must be
learned from incomplete observations, by approximately solving
\begin{equation}
\min_{D,\,A}\;
\sum_{q} \norm{M_q\,(R_q y - D\alphavec_q)}_2^2
\quad\text{s.t.}\;\norm{\alphavec_q}_0\le s,\;\norm{D_{:,k}}_2=1,
\label{eq:masked-loss}
\end{equation}
where $M_q\in\{0,1\}^{n\times n}$ is the diagonal restriction of $M$
to the patch at position $q$. Equation~\eqref{eq:masked-loss} is the
established \emph{weighted K-SVD objective}, due to
Mairal--Elad--Sapiro~\cite{mairal2008sparse} and refined in
Naumova--Schnass~\cite{naumova2018itkrmm} and
Tiefenthaler~\cite{tiefenthaler2018thesis}; it is \emph{not} our
contribution.

The remainder of this proposal addresses the practical and theoretical
\emph{gap} between existing solvers of \eqref{eq:masked-loss} and what
is required for high-fidelity inpainting under (a) heavy missingness
($r>70\%$), (b) holes that are larger than a patch, and (c) mixed
smooth/textured content -- regimes where every existing single-scale,
binary-masked, non-coupled solver in the literature is observably
suboptimal.

% =====================================================================
\section{Related Work and the State of the Art}
% =====================================================================
\label{sec:related}

\subsection{Masked dictionary inpainting (the baseline)}

The objective \eqref{eq:masked-loss} originates with
Mairal--Elad--Sapiro~\cite{mairal2008sparse} (\emph{weighted K-SVD},
``wKSVD''), where the dictionary update is performed by a weighted
rank-one approximation of the masked error matrix. Convergence of the
weighted rank-one step is established by
Srebro--Jaakkola~\cite{srebro2003weighted} via an EM/imputation argument.

The \emph{ITKrMM} algorithm of Naumova--Schnass~\cite{naumova2018itkrmm}
extends this line by replacing OMP with thresholding and the SVD with
$K$ residual means, and additionally learns an \emph{$L$-dimensional
low-rank component} $\Gamma$ that captures the dense ``mean +
contrast'' part of natural-image statistics.

Most recently, Tiefenthaler~\cite{tiefenthaler2018thesis} (Innsbruck,
2018, supervisor K.~Schnass) explicitly extends wKSVD to allow
arbitrary-dimensional low-rank components (algorithm \emph{wKSVDL}),
and benchmarks wKSVDL against ITKrMM on Peppers and Mandrill at
$30\%/50\%/70\%$ random erasure plus a structured \emph{Cracks} mask.
The thesis settles, in the Schnass group's own line of work, the
\emph{single-scale}, binary-mask, single-pass dictionary inpainting
question.

\textbf{What this prior art establishes (and our project does not
claim as novel).}\quad
The following are settled and we use them as baselines:
\begin{enumerate}[label=(B\arabic*),leftmargin=2.5em,itemsep=2pt]
\item \emph{Masked OMP} for sparse coding of partially observed
  patches (\cite{mairal2008sparse}, also Alg.~1
  of~\cite{tiefenthaler2018thesis}).
\item \emph{Weighted rank-one dictionary update} solving
  \eqref{eq:masked-loss}'s atom-wise sub-problem
  (\cite{srebro2003weighted,mairal2008sparse,tiefenthaler2018thesis}).
\item \emph{Single-pass dictionary inpainting} -- learn $D$ once,
  sparse-code masked patches, overlap-average to a final image.
\item \emph{Low-rank component handling} (ITKrMM and wKSVDL).
\end{enumerate}

\subsection{Smooth-prior inpainting}

Variational and PDE-based methods (Chan--Shen
\cite{chan2002mathematical}, Bertalmio--Sapiro
\cite{bertalmio2000image}) provably dominate dictionary-based methods
on smooth holes: biharmonic completion is the unique minimiser of
bending energy with Dirichlet/Neumann boundary conditions, hence
essentially optimal in the texture-free regime. Conversely, these
methods smear textures.

\subsection{Non-local sparse representations}

Centralised Sparse Representation (CSR,
Dong--Zhang--Shi~\cite{dong2011centralized}) and Group-Based Sparse
Representation (GSR,
Zhang--Zhao--Gao~\cite{zhang2014groupbased}) showed that regularising
each patch's sparse code toward the weighted mean of its similar
patches' codes improves stability on \emph{fully observed}
restoration. To our knowledge, no prior work integrates this kind of
non-local coefficient coupling into a \emph{masked, iterative}
dictionary-inpainting solver with provable error control. CROWN-Inpaint
does.

\subsection{Multi-scale image processing}

Image pyramids and coarse-to-fine optimisation are classical, but their
specific role in escaping the \emph{zero-information plateau} of
masked K-SVD when the hole exceeds the patch has not, to our
knowledge, been formally analysed in the dictionary-inpainting
literature.

% =====================================================================
\section{The Gap We Exploit}
% =====================================================================
\label{sec:gap}

The four limitations of every solver of \eqref{eq:masked-loss} known
to us, including~\cite{mairal2008sparse,naumova2018itkrmm,tiefenthaler2018thesis},
are the following.

\paragraph{L1: binary masks.}
All cited works treat the mask as $M\in\{0,1\}^d$. A pixel inside the
hole is either ``trusted at $0$'' or ``not used''. There is no
mechanism for the solver to assign \emph{partial} trust to a hole
estimate and update that trust over outer iterations.

\paragraph{L2: no smooth/sparse fusion.}
Pure dictionary inpainting smears smooth regions; pure biharmonic
inpainting smears textures. Existing solvers commit to one prior. No
existing masked-dictionary solver fuses both with a per-pixel,
data-driven fusion weight.

\paragraph{L3: no non-local coupling inside the masked solver.}
CSR and GSR demonstrated the value of non-local coefficient coupling on
fully observed data, but not within a confidence-weighted, masked,
iterative framework, and not with explicit error bounds.

\paragraph{L4: no remedy for holes $>$ patch size.}
When a contiguous hole is larger than the patch, every patch interior
to the hole has $M_q=\bm{0}$, forcing $\alphavec_q=\bm{0}$ at every
iteration -- a \emph{zero-information plateau} on the loss surface.
Prior work either avoids this regime experimentally or accepts the
resulting black smear.

\medskip\noindent
\textbf{CROWN-Inpaint addresses L1--L4 simultaneously, with formal
guarantees, while keeping the masked-K-SVD backbone of B1--B4 intact.}

% =====================================================================
\section{Proposed Contribution: CROWN-Inpaint}
% =====================================================================
\label{sec:proposal}

\subsection{High-level architecture}

CROWN-Inpaint is an \emph{outer iterative} solver. Each outer iteration
performs five steps:

\begin{enumerate}[label=\textbf{Step \arabic*.},leftmargin=4.2em,itemsep=3pt]
\item \textbf{Confidence map update.} Compute a continuous $c^t\in[0,1]^{H\times W}$ on
hole pixels from (a) overlap variance of patch reconstructions, (b)
distance to $\Omega$, and (c) an iteration schedule. Combine with the
binary mask to form $W^t = M+(1-M)\odot c^t$.
\item \textbf{Confidence-weighted sparse branch.} For every patch
$q$, run a weighted-OMP on $(R_q u^{t-1}, D, W^t_q)$, then refine the
selected codes by a non-local centralised $\ell_1$ ISTA toward the
group centre $\bm{\beta}_q$. Overlap-average to obtain
$u^t_{\rm sparse}$.
\item \textbf{Smooth branch.} Run a small number of harmonic Jacobi
relaxation steps from $u^{t-1}$, holding $\Omega$ fixed, to obtain
$u^t_{\rm smooth}$.
\item \textbf{Regime-aware fusion.} Compute or update a regime map
$r\in[0,1]^{H\times W}$ from observed gradients, structure-tensor
anisotropy and DCT spectral entropy; fuse
$\bar u^t = (1-r)\odot u^t_{\rm smooth} + r\odot u^t_{\rm sparse}$.
\item \textbf{Hard projection.}
$u^t = M\odot y + (1-M)\odot \bar u^t$. Optionally apply a small
stochastic manifold-correction step every few iterations.
\end{enumerate}

The outer loop is wrapped inside a \emph{three-level coarse-to-fine
pyramid} that performs the entire procedure at $32\times 32$, then
$64\times 64$, then $128\times 128$, with the coarse reconstruction
injected as a warm start to the next finer level.

\subsection{The four genuinely novel mechanisms}

We list the four \emph{additive} novel mechanisms that CROWN-Inpaint
introduces above the baselines B1--B4 and the gaps L1--L4.

\paragraph{N1 -- Continuous confidence-weighted OMP and outer loop.}
We replace the binary mask $M$ used by every prior masked-dictionary
solver by a continuous confidence-weighted $W^t\in[0,1]^d$. Concretely
we solve, per patch and per outer iteration,
\begin{equation}
\min_{\alphavec}\;\norm{(W^t_q)^{1/2}\odot(R_q u^{t-1}-D\alphavec)}_2^2
\quad\text{s.t.}\;\norm{\alphavec}_0\le s.
\label{eq:weighted-omp-prob}
\end{equation}
The substitution $\widetilde D = \diag\bigl((W^t_q)^{1/2}\bigr)D$,
$\widetilde x = (W^t_q)^{1/2}\odot R_q u^{t-1}$ reduces
\eqref{eq:weighted-omp-prob} \emph{exactly} to standard OMP on
$(\widetilde D,\widetilde x)$ with per-column re-normalisation,
preserving the reduction trick of binary masked OMP while
strictly generalising it. The confidence schedule is the key
mechanism that lets the solver gradually trust its own hole estimates
without admitting them blindly at iteration~1.

\paragraph{N2 -- Regime-aware fusion of smooth and sparse priors.}
We introduce a per-pixel regime map $r\in[0,1]^{H\times W}$ computed
on the \emph{observed} pixels via gradient energy, structure-tensor
anisotropy, and DCT spectral entropy, and propagated into the hole by
boundary-conditional Gaussian smoothing. The fusion
\[
\bar u = (1-r)\odot u_{\rm smooth} + r\odot u_{\rm sparse}
\]
gates two priors with provably complementary regimes of optimality
(biharmonic on smooth holes; sparse coding on textured holes). To our
knowledge no prior masked-dictionary inpainting work performs this
fusion at all, let alone with a learned per-pixel gate.

\paragraph{N3 -- Non-local centralised $\bm{\ell_1}$ coupling inside
the masked solver.}
For each patch $q$, define a non-local prior
\begin{equation}
\bm{\beta}_q = \sum_{r\in\KNN(q)} w_{qr}\,\alphavec_r,\qquad
w_{qr}=\frac{\exp(-\norm{R_q u-R_r u}_{W_q\odot W_r}^2/h^2)}{Z_q},
\end{equation}
where the similarity weights use \emph{coordinates confidently
observed in both patches}. Then refine each code by
\begin{equation}
\alphavec_q^\star
= \arg\min_{\alphavec}\;\norm{W_q^{1/2}(R_q u - D\alphavec)}_2^2
+ \lambda\norm{\alphavec - \bm{\beta}_q}_1.
\label{eq:nsr}
\end{equation}
This is a \emph{convex} sub-problem solved by an ISTA loop of $\sim 20$
iterations on the OMP support. Inspired by CSR
\cite{dong2011centralized} and GSR \cite{zhang2014groupbased},
\eqref{eq:nsr} differs in three concrete ways: the data term is
\emph{confidence-weighted}, the prior $\bm{\beta}_q$ uses
\emph{double-masked} similarity (so unobserved coordinates do not
spuriously influence neighbour selection), and the refinement is
\emph{embedded inside} an iterative masked solver rather than applied
once on fully observed data.

\paragraph{N4 -- Coarse-to-fine pyramid with prior injection and
plateau-escape guarantee.}
Existing single-scale solvers fail catastrophically when the hole side
$H_\Omega$ exceeds the patch side $p$, because every patch interior to
the hole has $M_q=\bm{0}$, leading to $\nabla_{\alphavec_q}\Loss\equiv\bm{0}$
and an irreducible zero-information plateau on the loss surface.
CROWN-Inpaint runs the entire algorithm at the coarsest pyramid level
where $H_\Omega/2^{L-1}\le p$ -- at which the dead set is empty -- and
uses bilinear-upsampled coarse reconstructions as warm starts at
finer levels. We \emph{prove} (Theorem~\ref{thm:warm} below) that the
warm start strictly removes the plateau and gives the descent
algorithm a non-zero gradient on which to act.

\subsection{What CROWN-Inpaint does not claim as novel}

For complete intellectual honesty:
\begin{enumerate}[label=$\circ$,leftmargin=2.4em,itemsep=2pt]
\item Masked OMP, masked K-SVD, and the weighted rank-one dictionary
  update are baselines from
  \cite{mairal2008sparse,srebro2003weighted,naumova2018itkrmm,tiefenthaler2018thesis}.
\item The closed-form scalar coordinate ALS we use to implement the
  rank-one update is a \emph{coordinate-descent variant} of the
  EM/SVD weighted rank-one update of~\cite{srebro2003weighted}; we
  re-derive it because the scalar form is faster and avoids any
  imputation step, but the underlying objective is identical.
\item The biharmonic smooth branch is~\cite{chan2002mathematical}
  used as-is.
\item The non-local idea of regularising codes toward neighbour means
  is from CSR/GSR; what is novel is its concrete integration into
  the masked, iterative, confidence-weighted setting (item N3 above).
\item The pyramid is a classical idea; what is novel is the formal
  plateau-escape guarantee in the masked-K-SVD setting (item N4).
\item Plug-and-play stochastic manifold correction is from the
  diffusion-PnP literature; we use it as an \emph{optional} stabiliser.
\end{enumerate}

% =====================================================================
\section{Mathematical Formulation and Inspiration}
% =====================================================================

We can view CROWN-Inpaint as approximate alternating minimisation of
the energy
\begin{align}
E(u,\{\alphavec_q\})
\;=\;& \sum_q \norm{W_q^{1/2}\odot(R_q u - D\alphavec_q)}_2^2
       \;+\; \lambda \sum_{q,\,r\in\KNN(q)} w_{qr}\norm{\alphavec_q-\alphavec_r}_1
       \nonumber\\
     &\;+\; \mu\, J_{\rm smooth}(u;\,r)
       \;+\; \iota_\Omega(u;\,y),
\label{eq:energy}
\end{align}
subject to $\norm{\alphavec_q}_0\le s$. The four terms correspond
exactly to the four CROWN priors:
\begin{itemize}[itemsep=2pt]
\item the \emph{confidence-weighted} data term (mechanism N1) reduces to
classical masked sparse coding when $W_q=M_q\in\{0,1\}^d$;
\item the \emph{non-local} term (N3) corresponds to a convex relaxation
of the $\ell_1$ deviation from $\bm{\beta}_q$ when one notes that
$\sum_r w_{qr}\norm{\alphavec_q-\alphavec_r}_1 \approx
\norm{\alphavec_q-\bm{\beta}_q}_1$ to leading order in the centred
shrinkage analysis;
\item the \emph{smooth} term (N2) is a Tikhonov / harmonic energy
locally modulated by the regime map $r$;
\item the indicator $\iota_\Omega$ implements the hard observation
constraint (Step~5).
\end{itemize}
Each step of the outer iteration corresponds to one block-coordinate
descent step on \eqref{eq:energy} with a closed-form (or
near-closed-form) inner solve. The pyramid (N4) corresponds to
warm-starting the non-convex problem from a coarse minimiser.

\subsection{Inspirations -- one line each}

\begin{itemize}[leftmargin=*,itemsep=2pt]
\item \textbf{Patch sparsity}: K-SVD~\cite{aharon2006ksvd}, masked
extensions~\cite{mairal2008sparse,naumova2018itkrmm,tiefenthaler2018thesis}.
\item \textbf{Confidence weighting}: Plug-and-play and reliability-weighted
estimation, recast for masked OMP.
\item \textbf{Regime / fusion}: classical observation that biharmonic is
optimal on smooth holes \cite{chan2002mathematical} and that texture
needs atom transport from the boundary~\cite{bertalmio2000image}.
\item \textbf{Non-local centralised refinement}: CSR
\cite{dong2011centralized}, GSR~\cite{zhang2014groupbased}, with the
slow-rate Lasso analysis of
Bickel--Ritov--Tsybakov~\cite{bickel2009simultaneous}.
\item \textbf{Multi-scale warm start}: the empirical observation that
single-scale K-SVD fails catastrophically when
$H_\Omega>p$, formalised via the plateau-escape theorem of
Sec.~\ref{sec:why-works}.
\item \textbf{Optional manifold correction}: stochastic plug-and-play
without requiring a pre-trained denoiser network.
\end{itemize}

% =====================================================================
\section{Why It Works: Theoretical Foundations}
% =====================================================================
\label{sec:why-works}

We summarise four results that we have already proven for the
companion convergence document; full proofs are in the technical
report.

\paragraph{T1: Convergence of the masked dictionary update.}
For each atom $k$, the per-atom rank-one sub-problem
\[
\min_{\dvec_k,\hvec_k}\sum_{i\in\Sk}\norm{M_i(\bm{e}_{k,i}-\dvec_k h_{k,i})}_2^2
\]
admits the closed-form scalar updates
\[
h_{k,i}=\frac{\sum_j (M_i)_{jj}\,d_{k,j}\,e_{k,i,j}}{\sum_j (M_i)_{jj}\,d_{k,j}^2},
\quad
d_{k,j}=\frac{\sum_i (M_i)_{jj}\,h_{k,i}\,e_{k,i,j}}{\sum_i (M_i)_{jj}\,h_{k,i}^2},
\]
each of which globally minimises the masked Frobenius cost in its
single block. Under a standard atom-utilisation assumption, alternating
the two updates yields a monotonically non-increasing energy converging
to a coordinate-wise stationary point, and the outer K-SVD loop
inherits this descent property. \emph{This recovers the known
guarantee of~\cite{srebro2003weighted}, restated as a coordinate
descent for clarity and speed.}

\paragraph{T2: Slow-rate Lasso bound for the non-local refinement
(novel).}
Define the \textbf{Self-Similarity Index}
\[
\mathrm{SSI}_q := \norm{\alphavec_q^{\rm true}-\bm{\beta}_q}_1.
\]

\begin{theorem}[Self-similarity error bound]
\label{thm:lasso}
Under the noise model $R_q u^{\rm true} = D\alphavec_q^{\rm true}+\bm{\nu}_q$,
$\norm{W_q^{1/2}\bm{\nu}_q}_2\le\delta$, and $\lambda\ge
2\norm{\widetilde D_q^\top \widetilde{\bm{\nu}}_q}_\infty$, every
solution $\alphavec_q^\star$ of \eqref{eq:nsr} satisfies
\[
\norm{W_q^{1/2}\bigl(R_q u^{\rm true}-D\alphavec_q^\star\bigr)}_2
\;\le\;\delta+\sqrt{3\lambda\cdot\mathrm{SSI}_q}.
\]
\end{theorem}

This is the slow-rate Lasso inequality applied to the centred
regulariser $\norm{\alphavec-\bm{\beta}_q}_1$ and is, to our
knowledge, the first explicit error bound for a masked, iterative,
non-locally coupled sparse-coding step. The proof proceeds by the
change of variables $\bm{\gamma}=\alphavec-\bm{\beta}_q$, which
reduces \eqref{eq:nsr} to a standard Lasso whose oracle inequality is
then applied. The bound captures the expected behaviour: better
self-similar matches (smaller $\mathrm{SSI}_q$) yield tighter
reconstruction error.

\paragraph{T3: Plateau-escape via multi-scale warm start (novel).}

\begin{lemma}[Plateau gradient]
\label{lem:plateau}
For every patch $q$ with $M_q=\bm{0}$ and any $D$,
\[
\nabla_{\alphavec_q}\Loss(D,A)
\;=\;-2D^\top M_q(R_q u-D\alphavec_q)\;\equiv\;\bm{0}.
\]
\end{lemma}

\begin{theorem}[Warm-start removes the plateau]
\label{thm:warm}
Let $\widehat u^{(\ell)}$ be the reconstruction at the coarse scale,
$\mathcal{U}_{\ell\to\ell-1}$ the bilinear upsampler, and define the
prior-injected initialisation
\[
u^{(\ell-1)}_{\rm init}
=M^{(\ell-1)}\odot y^{(\ell-1)}+(1-M^{(\ell-1)})\odot \mathcal{U}_{\ell\to\ell-1}\bigl(\widehat u^{(\ell)}\bigr).
\]
If the upsampling error is at most $\eta$ on the hole and the coarse
scale has empty dead set, then the gradient of the prior-augmented
loss at $u^{(\ell-1)}_{\rm init}$ is generically non-zero, the OMP
sparse coding step returns codes with per-patch fitting error
$\le n\eta^2$, and one ALS sweep produces a strictly lower
prior-augmented loss. By contrast, any cold-started iterate
$(D^{(0)},A=\bm{0})$ on the strict-mask loss is permanently stuck:
$A^{(t)}_{:,q}=\bm{0}$ for all $t\ge 0$ and every $q$ with $M_q=\bm{0}$.
\end{theorem}

\paragraph{T4: Composite end-to-end guarantee.}
Combining T1--T3, the full CROWN-Inpaint pipeline produces a sequence
$(D^{(t)},A^{(t)})$ at each pyramid level that is monotonically
non-increasing in the corresponding masked energy, with bounded
per-patch reconstruction error
$\le \delta_q+\sqrt{3\lambda\cdot\mathrm{SSI}_q}$ at termination of
the non-local step. Crucially, T3 ensures that this descent is taken
\emph{from a non-trivial starting point} on holes larger than the
patch -- the regime where every prior solver fails.

\begin{remark}
None of T1--T4 claim global optimality of the recovered $(D,A)$:
$\ell_0$ sparse coding is NP-hard and OMP is a greedy approximation.
This caveat is identical to that of classical
K-SVD~\cite{aharon2006ksvd} and is explicitly inherited by every
masked-K-SVD-based method in the literature.
\end{remark}

% =====================================================================
\section{Experimental Plan}
% =====================================================================

\subsection{Datasets and masks}
\begin{description}[leftmargin=2.5em,itemsep=2pt]
\item[Datasets.] \emph{Peppers} and \emph{Mandrill} ($256\times 256$,
  for direct comparison with~\cite{tiefenthaler2018thesis}); subsets
  of \emph{CelebA-HQ} and \emph{Places2} ($128\times 128$ patches);
  textured subsets from DTD for stress tests.
\item[Masks.] (a) Random Bernoulli erasure at $30\%/50\%/70\%$
  (matches~\cite{tiefenthaler2018thesis}); (b) the same \emph{Cracks}
  structured mask used in~\cite{tiefenthaler2018thesis}; (c)
  rectangular block holes of side $20{-}40$ at $128\times 128$ to
  exercise the multi-scale plateau-escape mechanism (no prior solver
  handles this gracefully); (d) thin scratches and free-form brush
  masks.
\end{description}

\subsection{Baselines}
We benchmark against the following, each implemented or wrapped in our
codebase:
\begin{enumerate}[label=(\alph*),itemsep=2pt]
\item \textbf{Biharmonic / Telea inpainting} (classical, on smooth holes).
\item \textbf{Plain masked K-SVD} (B1+B2+B3 of Sec.~\ref{sec:related},
  our own implementation already complete).
\item \textbf{wKSVDL} of Tiefenthaler~\cite{tiefenthaler2018thesis} (we
  will re-implement Algorithms 4--5 of that thesis as our principal
  competing single-scale baseline).
\item \textbf{ITKrMM} of Naumova--Schnass~\cite{naumova2018itkrmm}.
\item \textbf{CSR/GSR-style} non-local sparse representation
  \cite{dong2011centralized,zhang2014groupbased}, where feasible
  given the masked setting.
\item \textbf{CROWN-Inpaint} (full method, our contribution).
\item Deep baselines (\emph{stretch goal}): a single LaMa or
  partial-conv pre-trained model, as a sanity reference -- we are an
  optimisation-first method and do not aim to beat large pretrained
  models.
\end{enumerate}

\subsection{Ablations (the experimental core of the paper)}

We will run ablation rows that toggle each of the four CROWN mechanisms
on/off independently:

\begin{center}
\begin{tabular}{@{}lcccc@{}}
\toprule
Variant & N1 conf. & N2 regime & N3 nonlocal & N4 pyramid \\
\midrule
Plain masked K-SVD & -- & -- & -- & -- \\
+N1 only           & \checkmark & -- & -- & -- \\
+N2 only           & -- & \checkmark & -- & -- \\
+N3 only           & -- & -- & \checkmark & -- \\
+N4 only           & -- & -- & -- & \checkmark \\
+N1+N3             & \checkmark & -- & \checkmark & -- \\
+N1+N4             & \checkmark & -- & -- & \checkmark \\
\textbf{CROWN-Inpaint (full)} & \checkmark & \checkmark & \checkmark & \checkmark \\
\bottomrule
\end{tabular}
\end{center}

This table is precisely what reviewers will demand and exactly what
distinguishes us from~\cite{tiefenthaler2018thesis}, which only varies
the low-rank dimension $L\in\{1,3,8\}$.

\subsection{Metrics}
PSNR, SSIM, hole-only PSNR, boundary MAE, mean wall-clock per image,
and number of outer iterations to convergence. We will report mean and
standard deviation over five random seeds per (image, mask) pair, as
in~\cite{tiefenthaler2018thesis}.

\subsection{Hypotheses we will test}

\begin{description}[leftmargin=2.4em,itemsep=2pt]
\item[H1.] N1 (continuous confidence) gives a $\ge 0.5$ dB
  hole-PSNR improvement at $r\ge 50\%$ over binary-mask wKSVDL.
\item[H2.] N2 (regime fusion) eliminates the texture-vs-smooth
  trade-off observed by~\cite{tiefenthaler2018thesis}: the same
  CROWN configuration wins both Peppers (smooth) and Mandrill
  (textured), unlike wKSVDL which loses to ITKrMM on Mandrill.
\item[H3.] N3 (non-local) reduces variance of reconstruction across
  random seeds by $\ge 30\%$ at $r=70\%$.
\item[H4.] N4 (pyramid) recovers semantically meaningful content in
  $\ge 20\times 20$ block holes where every single-scale baseline
  produces a black smear.
\end{description}

% =====================================================================
\section{Implementation Status}
% =====================================================================

The codebase is already organised around the architecture above.

\begin{center}
\small
\begin{tabular}{@{}lll@{}}
\toprule
File & Component & Status \\
\midrule
\texttt{masked\_ksvd.py}              & Masked OMP + ALS update (B1, B2)        & complete \\
\texttt{inpainting\_masked\_ksvd.py}  & Single-scale masked inpainting           & complete \\
\texttt{inpainting\_multiscale\_masked\_ksvd.py} & 3-level pyramid (N4)        & complete \\
\texttt{crown/confidence.py}          & Confidence map (N1)                      & complete \\
\texttt{crown/regime.py}              & Regime map (N2)                          & complete \\
\texttt{crown/weighted\_omp.py}       & Confidence-weighted OMP (N1)             & complete \\
\texttt{crown/nonlocal\_coupling.py}  & Non-local centralised ISTA (N3)          & complete \\
\texttt{crown/smooth.py}              & Biharmonic / harmonic Jacobi             & complete \\
\texttt{crown/fuse.py}                & Regime-aware fusion + projection         & complete \\
\texttt{crown/manifold.py}            & Optional stochastic correction           & complete \\
\texttt{crown/run.py}                 & Outer-loop driver                        & complete \\
\texttt{experiment\_sweep.py}         & Ablation grid                            & in progress \\
wKSVDL re-implementation              & Tiefenthaler 2018 baseline              & to do \\
\bottomrule
\end{tabular}
\end{center}

The convergence-and-error analysis is also written up as a separate
technical companion document \texttt{convergence\_proofs.tex}
(IEEE format, 9 pages). It contains full proofs of T1--T4.

% =====================================================================
\section{Timeline and Deliverables}
% =====================================================================

\begin{center}
\small
\begin{tabular}{@{}llp{8.5cm}@{}}
\toprule
Week & Milestone & Notes \\
\midrule
W1   & wKSVDL re-implementation & Algorithms 4--5 of~\cite{tiefenthaler2018thesis}, validated to match the thesis's Peppers/Mandrill PSNRs ($\pm 0.3$ dB). \\
W2   & Ablation sweep          & 8-row ablation $\times$ 4 mask types $\times$ 2 images $\times$ 5 seeds. \\
W3   & Multi-scale stress test & $20\times 20$ and $40\times 40$ block holes; visual + numerical comparison with all single-scale baselines. \\
W4   & Paper draft -- main     & Sections 1--4 (introduction, related work, method, theory). \\
W5   & Paper draft -- experiments & Section 5; integrate plots from W2--W3. \\
W6   & Internal review         & Course staff feedback; revise. \\
W7   & Final submission        & PDF + reproducible code repository + raw experimental data. \\
\bottomrule
\end{tabular}
\end{center}

\noindent\textbf{Final deliverables:}
\begin{enumerate}[label=\textbf{D\arabic*.},leftmargin=2.6em,itemsep=2pt]
\item Paper PDF (8--10 pages, IEEE format) describing the method,
  theory, and experiments.
\item Companion convergence-proof document
  (\texttt{convergence\_proofs.tex}, already drafted).
\item Reproducible code repository with all eight ablation rows and
  every baseline.
\item Experimental data (\texttt{.csv} + \texttt{.npy}) and a
  notebook reproducing every figure in the paper.
\end{enumerate}

% =====================================================================
\section{Risks and Mitigations}
% =====================================================================

\paragraph{R1: Tiefenthaler 2018 covers more of the masked-dictionary
ground than expected.}
\emph{Already incurred and accommodated.} The present proposal
explicitly demotes masked K-SVD itself to a baseline and re-centres
novelty on the four CROWN mechanisms (N1--N4). No further pivot is
required.

\paragraph{R2: Non-local refinement is too slow for $128\times 128$
images.}
We restrict the centralised $\ell_1$ refinement to the OMP support
(at most $s\le 12$ atoms), turning a $K$-dim convex problem into an
$s$-dim one. ISTA on $s$-dim problems is $O(s^2)$ per iteration with
$\sim 20$ iterations; total cost is dominated by the OMP step.

\paragraph{R3: The pyramid leaks coarse-scale artefacts upward.}
We perform mask-aware downsampling (normalised convolution that
excludes hole pixels from the aggregation) at every level transition;
this is already implemented in
\texttt{inpainting\_multiscale\_masked\_ksvd.py}. We will additionally
ablate the upsampling kernel.

\paragraph{R4: Manifold correction destabilises convergence.}
We treat it as optional (off by default) and include it only as an
ablation knob, with a decreasing $\sigma_t=\sigma_0\gamma^{t-1}$
schedule that bounds its long-run effect.

\paragraph{R5: Reviewers question whether N1+N2+N3+N4 are independent
or redundant.}
The 8-row ablation table directly answers this: each row toggles a
single mechanism, and we will report the marginal PSNR contribution of
each.

% =====================================================================
\section{Expected Outcomes and Publication Target}
% =====================================================================

\begin{itemize}[itemsep=3pt]
\item A complete, theoretically-grounded, optimisation-first
  inpainting method that beats every single-scale masked-dictionary
  baseline (including wKSVDL) by $1{-}3$ dB on standard benchmarks
  and produces visually coherent fills on holes that are
  $2{-}5\times$ the patch size.
\item Four formal results (T1--T4) with full proofs; T2--T4 are, to
  our knowledge, new in the masked-dictionary inpainting literature.
\item A clean ablation table that pinpoints the marginal contribution
  of each CROWN mechanism.
\item A reproducible code repository.
\end{itemize}

\noindent\textbf{Submission target.} A workshop or short paper at
ICIP / ICASSP 2026, or a journal extension to \emph{IEEE Trans.\ Image
Processing} or \emph{IEEE Trans.\ Computational Imaging}; we will
make a final venue decision after the W6 internal review based on
experimental strength.

% =====================================================================
\section*{Acknowledgements}
% =====================================================================
We thank the course staff of CS 754 (Spring 2026, IIT Bombay) for
supervision and feedback, and explicitly acknowledge
M.~Tiefenthaler and K.~Schnass~\cite{tiefenthaler2018thesis} for the
clearest exposition of single-scale masked K-SVD inpainting we have
encountered, against which CROWN-Inpaint is benchmarked.

% =====================================================================
\begin{thebibliography}{99}
% =====================================================================

\bibitem{aharon2006ksvd}
M.~Aharon, M.~Elad, and A.~Bruckstein,
``K-SVD: An algorithm for designing overcomplete dictionaries for sparse
representation,''
\emph{IEEE Trans. Signal Process.}, vol.~54, no.~11, pp.~4311--4322, 2006.

\bibitem{mairal2008sparse}
J.~Mairal, M.~Elad, and G.~Sapiro,
``Sparse representation for color image restoration,''
\emph{IEEE Trans. Image Process.}, vol.~17, no.~1, pp.~53--69, 2008.

\bibitem{srebro2003weighted}
N.~Srebro and T.~Jaakkola,
``Weighted low-rank approximations,''
in \emph{Proc.\ Int.\ Conf.\ Machine Learning (ICML)}, 2003,
pp.~720--727.

\bibitem{naumova2018itkrmm}
V.~Naumova and K.~Schnass,
``Dictionary learning from incomplete data for efficient image
restoration,''
in \emph{Proc.\ European Signal Processing Conf.\ (EUSIPCO)}, 2017
(see also their full ITKrMM manuscript, \emph{Inverse Problems and
Imaging}, 2018).

\bibitem{tiefenthaler2018thesis}
M.~Tiefenthaler,
``Integrating low-rank components into weighted K-SVD for dictionary
based inpainting,''
Bachelor thesis, Department of Mathematics, University of Innsbruck,
Austria, supervised by K.~Schnass, September 2018.

\bibitem{tropp2007omp}
J.~A.~Tropp and A.~C.~Gilbert,
``Signal recovery from random measurements via orthogonal matching
pursuit,''
\emph{IEEE Trans. Inf. Theory}, vol.~53, no.~12, pp.~4655--4666, 2007.

\bibitem{dong2011centralized}
W.~Dong, L.~Zhang, and G.~Shi,
``Centralized sparse representation for image restoration,''
in \emph{Proc.\ IEEE Int.\ Conf.\ Computer Vision (ICCV)}, 2011,
pp.~1259--1266.

\bibitem{zhang2014groupbased}
J.~Zhang, D.~Zhao, and W.~Gao,
``Group-based sparse representation for image restoration,''
\emph{IEEE Trans. Image Process.}, vol.~23, no.~8, pp.~3336--3351,
2014.

\bibitem{chan2002mathematical}
T.~F.~Chan and J.~Shen,
``Mathematical models for local nontexture inpaintings,''
\emph{SIAM J. Appl. Math.}, vol.~62, no.~3, pp.~1019--1043, 2002.

\bibitem{bertalmio2000image}
M.~Bertalmio, G.~Sapiro, V.~Caselles, and C.~Ballester,
``Image inpainting,''
in \emph{Proc.\ ACM SIGGRAPH}, 2000, pp.~417--424.

\bibitem{bickel2009simultaneous}
P.~J.~Bickel, Y.~Ritov, and A.~B.~Tsybakov,
``Simultaneous analysis of Lasso and Dantzig selector,''
\emph{Ann. Statist.}, vol.~37, no.~4, pp.~1705--1732, 2009.

\bibitem{beck2009fista}
A.~Beck and M.~Teboulle,
``A fast iterative shrinkage-thresholding algorithm for linear inverse
problems,''
\emph{SIAM J. Imaging Sci.}, vol.~2, no.~1, pp.~183--202, 2009.

\bibitem{tseng2001convergence}
P.~Tseng,
``Convergence of a block coordinate descent method for nondifferentiable
minimization,''
\emph{J. Optim. Theory Appl.}, vol.~109, no.~3, pp.~475--494, 2001.

\end{thebibliography}

\end{document}
